% !TeX program = xelatex
% 第二研究步骤独立任务报告；在项目根目录运行 make intermediate。
\documentclass[UTF8,a4paper,11pt,oneside,fontset=fandol]{ctexart}
\newcommand{\RunningTitle}{第二研究步骤：日度流动性与可交割供给}
\input{intermediate/templates/shared.tex}
\usepackage{amsmath}
\usepackage{xurl}
\begin{document}
\thispagestyle{plain}
\begin{center}
{\LARGE\bfseries 第二研究步骤任务报告：\par}
{\LARGE\bfseries CEA、CCER 日度流动性与合格资产供给\par}
\vspace{0.4em}
{\small 核验日期：2026 年 9 月 30 日}\par
\end{center}

\section{任务范围和判断}
本步骤依课程合约设计要求，取得同一观察期的 CEA 与 CCER 逐日价量，按交易方式检查成交连续性、集中度和价格变化，并核实余额与账户权限能否支持拟议期货交割。观察样本为 2025 年全国市场现货，不含地方试点配额。原始行、来源链接及提取程序位于\par\noindent\path{research/contract-design/data/}\par\noindent 和 \path{research/contract-design/build_2025_daily_data.py}；完整核查底稿见\par\noindent\path{research/contract-design/DAILY_LIQUIDITY_AND_DELIVERABLE_SUPPLY_2025.md}。

\textbf{结论：继续把全国 CEA 作为研究性暂选标的。}在同一 203 日窗口里，CEA 挂牌市场的量明显高于 CCER，且合规清缴形成明确的潜在套保需求。但现有公开资料尚不能给出指定年度、未冻结、可划转的 CEA 余额，也未证明期货交割账户与登记结算接口已开通；2025 年综合价不能直接代表未来指定年度配额。因此本报告不确定交割年度、交割方式或合约数值参数。

\section{来源、清洗和年报对账}
CEA 来自上海环境能源交易所 2025 年每日概况目录，共 243 条，逐条保存公告 URL、综合收盘价、挂牌协议、大宗协议、单向竞价及总量额（R1）。遇公告明确写“无成交”，将量额记零，仍保留公告给出的综合收盘价，并标出非成交价日。交易方式未显示时，仅在已报分方式量额之和等于当日总量额时补零。243 条的量额及三种方式合计，逐吨逐分对上交易所全年汇总（R2）。工作区新提供的 CSMAR 日表中 2025 年全国 CEA 有 233 条成交日记录；其总量、总额、收盘价及三种交易方式的量额，九项字段逐日与官方提取表完全一致。CSMAR 表省略的恰为官方十个零成交日，且未细分配额年度；本地文件没有可核实的下载版本元数据，故只作交叉核查，不替代官方公告（R12）。

CCER 官方每日行情和年度一览为第一手来源（R3）；但官方日行情目录对本次程序访问返回 403，无法完整批量提取。逐日分析遂使用标注依据北京绿色交易所行情整理的 203 行公开转载表（R4），并以官方单日公告和累计数逐段对账。该表的 2025 年 6 月 24 日成交额应为 54,720.20 元，而非 54,720.00 元；7 月 1 日成交量应为 601 吨，而非 610 吨，已按官方公告更正（R5）。

\begin{table}[htbp]
\centering\small
\caption{2025 年量额核对及 CCER 转载表差额}
\begin{tabularx}{\textwidth}{@{}L{0.43\textwidth}XX@{}}
\toprule
口径 & 成交量（吨） & 成交额（元）\\
\midrule
CEA 官方逐日求和及年报 & 234,597,856 & 14,629,891,128.06\\
CCER 已作两条官方更正的转载表 & 8,835,094 & 625,169,946.60\\
CCER 官方累计数推得 2025 年全年 & 8,844,114 & 625,801,436.80\\
CCER 仍未定位至具体日的差数 & 9,020 & 631,490.20\\
\bottomrule
\end{tabularx}
\end{table}

CCER 精确全年量额的计算为：2025 年末官方累计 9,219,429 吨、649,636,716.80 元（12 月 30 日数加 12 月 31 日 1 吨、85 元），减去 2024 年首日 375,315 吨、23,835,280 元；2025 年 3 月 7 日官方累计数也核实了这个期初基数（R3、R6）。结果与官方年报的 884.41 万吨、6.26 亿元四舍五入口径相符。更正后的转载表在 9 月 10 日与官方累计数一致；至 9 月 15 日，官方累计比表内多 9,020 吨和 631,490.20 元（R7）。9 月 15 日单日与转载表一致，疑点集中于 9 月 11—12 日，尚无可靠原始公告可将差额分配到具体日。\textbf{以下 CCER 日度分布为近似值，全年精确总量以官方累计数为准。}

\section{同日期、同挂牌方式的比较}
两个市场从 2025 年 3 月 7 日至 12 月 31 日恰有相同的 203 个观察日期。CEA 用\textbf{挂牌协议}量与 CCER 2025 年仅有的挂牌协议量比较；CEA 大宗和竞价不计入此栏。分位数采用最近秩法，零量交易日保留。结果如下：

\begin{table}[htbp]
\centering\small
\caption{共同 203 日窗口的挂牌协议日量}
\begin{tabularx}{\textwidth}{@{}L{0.49\textwidth}XX@{}}
\toprule
指标 & CEA & CCER 转载表\\
\midrule
量合计（吨） & 79,011,724 & 8,835,094\\
日量中位数（吨） & 288,406 & 2,700\\
日量第 90 百分位（吨） & 912,040 & 77,932\\
零挂牌成交日 & 3 & 表内 0\\
前 10 个高量日占窗口量 & 16.35\% & 65.40\%\\
第四季度占窗口量 & 58.93\% & 68.28\%\\
\bottomrule
\end{tabularx}
\end{table}

将 CEA 窗口挂牌量除以\textbf{官方} CCER 年量，$79{,}011{,}724/8{,}844{,}114\approx8.93$。这是实现成交量的对照，不能当作可交割供给倍数。CEA 同窗所有交易方式总量 231,100,418 吨，明显高于挂牌量；若未来用现货挂牌价作最终结算价，不能拿大宗量充当挂牌价格深度。CCER 转载表中有 72 日量不超过 1,000 吨，前 10 日占约三分之二；小额日均价易受单笔和项目组合影响。两市场第四季度交易集中；官方 CCER 年报的四季度占比 68.2\% 与转载表 68.28\% 基本相符（R3），但这仍不能确定某一种可交割项目的季节性供需。

CEA 综合收盘价的逐日范围为 51.24—97.01 元/吨；CCER 转载表的日成交均价范围为 50.84—107.36 元/吨。仅作\textbf{描述}，将 CEA 有挂牌成交的 229 个日价按相邻有效观测计算 $|P_t/P_{t-1}-1|$，其中位数约 0.58\%；CCER 203 个日均价同法计算约 3.02\%。前者是综合收盘价，后者是跨项目成交均价，且一吨级成交日也进入 CCER 样本；二者不是同质、可用于直接设定保证金或涨跌停板的收益率波动率。CEA 2025 年综合价亦不是 2026 年初才上线的 CEA25 年度品种历史价（R8）。

\section{可交割余额和权限核查}
生态环境部 2025 年度信息称：CCER 年末累计登记减排量 1,776.37 万吨，注册登记系统累计开户 6,106 家（R9）。登记总量并未扣除已注销、用途限制或非拟议品级的余额；账户数量也不等于有期货交割权限的账户数量。生态环境部 2025 年答复称，截至当年 6 月底 CEA 累计发放配额约 279 亿吨，并指出全国市场初期现货交易仅在发电行业重点排放单位之间开展，其他主体将稳妥推进参与（R10）。累计发放量跨年度且含已清缴、冻结或转让部分，不能视作可交割供给。

2019—2024 年 CEA 结转须符合资格和数量要求，申请后待结转配额冻结，回收后发放等量 2025 年度配额（R11）。CEA25 于 2026 年 1 月 5 日上线，交易所同时发布 2026 年综合价格编制方案（R8）。这些事实要求交割品级按年度和登记状态界定，却不提供指定年度可用余额。公开材料仍未证明 CEA/CCER 登记系统可按广期所期货到期流程冻结、划转并接收资产。故\textbf{可持有、可现货交易与可期货交割应分别判断}；本轮无法确认实物交割能力或据此定每手数量、持仓限额。

\section{对后续合约设计的约束}
暂选全国 CEA 的依据是可核对的清缴需求和同日期挂牌交易规模。交割选择须先补三项直接证据：指定年度 CEA 的逐日挂牌价量；到期窗口内按年度、冻结状态及账户类别划分的可用余额；登记与期货结算之间的授权及划转操作说明。若转向现金交割，还需核准价格源授权、指定年度成分和无成交/异常成交备用规则。未取得这些资料前，报告中的合约月份、交易单位、涨跌停板、保证金及限仓应保持待定。

\section*{资料来源}
\small
\noindent R1. 上海环境能源交易所（2025）。《2025 年全国碳市场每日概况》目录及 243 条每日公告。\url{https://overview.cneeex.com/qgtpfqjy/mrgk/2025n/}\par
\noindent R2. 上海环境能源交易所（2025）。《全国碳市场每年综合价格行情及成交信息 20250102—20251231》。\url{https://overview.cneeex.com/c/2025-12-31/497010.shtml}\par
\noindent R3. 全国温室气体自愿减排交易系统（2025）。《2025 年全国温室气体自愿减排交易市场运行情况一览》。\url{https://www.ccer.com.cn/wcm/ccer/html/2307ptggc1/20260101/090113981.shtml}\par
\noindent R4. 碳中和网（2025）。《全国温室气体自愿减排项目核证自愿减排量（CCER）交易信息（2025 年）》；二次整理表。\url{https://www.ccn.ac.cn/carbon-market/ccer/ccerdate/219.html}\par
\noindent R5. 全国温室气体自愿减排交易系统（2025）。6 月 24 日与 7 月 1 日每日行情。\url{https://www.ccer.com.cn/wcm/ccer/html/2502lshq/20250624/170135985.shtml}\quad\url{https://www.ccer.com.cn/wcm/ccer/html/2502lshq/20250701/170330466.shtml}\par
\noindent R6. 全国温室气体自愿减排交易系统（2024—2025）。2024 年 1 月 22 日、2025 年 3 月 7 日与 12 月 30 日每日行情。\url{https://www.ccer.com.cn/wcm/ccer/html/2401jysj1/20250306/112037992.shtml}\quad\url{https://www.ccer.com.cn/wcm/ccer/html/2502lshq/20250307/174500101.shtml}\quad\url{https://www.ccer.com.cn/wcm/ccer/html/2502lshq/20251230/151804011.shtml}\par
\noindent R7. 全国温室气体自愿减排交易系统（2025）。9 月 10 日、9 月 15 日每日行情，生态中国网按原平台注明来源转载。\url{https://www.eco.gov.cn/news_info/62488.html}\quad\url{https://www.eco.gov.cn/news_info/62583.html}\par
\noindent R8. 上海环境能源交易所（2025）。《关于上线“碳排放配额25”及调整综合价格行情的通知》。\url{https://overview.cneeex.com/c/2025-12-29/496985.shtml}\par
\noindent R9. 生态环境部（2026）。《2025 年全国碳市场平稳有序运行》。\url{https://www.mee.gov.cn/ywgz/ydqhbh/wsqtkz/202601/t20260101_1139528.shtml}\par
\noindent R10. 生态环境部（2025）。《对十四届全国人大三次会议第 7512 号建议的答复》。\url{https://www.mee.gov.cn/xxgk2018/xxgk/xxgk13/202511/t20251114_1133700.html}\par
\noindent R11. 生态环境部（2024）。《2019—2024 年度碳排放配额结转方案》。\url{https://www.mee.gov.cn/xxgk2018/xxgk/xxgk06/202407/W020240702328244215540.pdf}\par
\noindent R12. CSMAR《中国碳排放权交易信息表（日）》，工作区所存 GF\_CEMISSRIGHTTRADE.csv 及同目录字段说明；文件时间 2026 年 9 月 30 日，下载版本元数据待核实。\par
\end{document}
